% --- Chapter: Modules ---
\section{Modules}

CRISP supports modular code organization with the \texttt{use} keyword. Currently, \texttt{use posix;} is the only built-in module.

\begin{tcolorbox}[colback=codebg, colframe=darkpurple, colbacktitle=darkpurple, title=\textbf{\textcolor{codegold}{Modules}}, fonttitle=\bfseries, sharp corners]
\begin{lstlisting}
# Load the POSIX module
use posix;

# POSIX functions live under a namespaced hash
let pid = posix["fork"]();
let fd = posix["open"]("/tmp/test.txt", O_WRONLY | O_CREAT, 0o644);
posix["write"](fd, "Hello!\n", 7);
posix["close"](fd);

# posix is a regular Hash — you can inspect it
for func in posix {
    say func;  # lists all POSIX function names
}
\end{lstlisting}
\end{tcolorbox}

\begin{tcolorbox}[colback=lightlavender, colframe=softvioletdark, title=\textbf{\textcolor{darkpurple}{Design Note: Module System}}, fonttitle=\bfseries, sharp corners]
The module system is under active development. Currently, only \texttt{use posix;} is supported. Future releases will extend this to user-defined modules (\texttt{use mymodule;}), import paths, and selective imports. The namespaced hash pattern (each module creating its own hash) will be the standard approach for all modules.
\end{tcolorbox}
